%%%%%%%%%%%%%%%%%%%%%%%%%%%%%%%%%%%%%%%%%%%%%%%%%%%%%%%%%%%%%%%%%%%%%%%%%%%%%
%  Préambule — identité typographique et graphique du rapport
%  Compilation : pdfLaTeX
%%%%%%%%%%%%%%%%%%%%%%%%%%%%%%%%%%%%%%%%%%%%%%%%%%%%%%%%%%%%%%%%%%%%%%%%%%%%%

% ----------------------------------------------------------------- encodage
\usepackage[T1]{fontenc}
\usepackage[utf8]{inputenc}
% Le français est la langue principale (dernière de la liste) ; l'anglais n'est
% chargé que pour la césure correcte de l'abstract.
\usepackage[english,french]{babel}
\frenchsetup{ThinSpaceInFrenchNumbers=true}

% -------------------------------------------------------------------- fontes
% Libertinus : empattements pour le corps, linéale assortie pour les titres et
% les figures. L'ensemble partage les mêmes proportions, ce qui évite la
% rupture visuelle habituelle entre texte et schémas.
\usepackage{libertinus}
\usepackage[varqu,varl]{inconsolata}   % monospace lisible pour le code
\usepackage{microtype}                 % césure et crénage optiques

% ------------------------------------------------------------------ couleurs
% Palette dérivée des logos de l'institut : le bleu est celui du logo INPT,
% le rouge celui du logo ANRT. Chaque couleur porte un sens unique et s'y
% tient dans tout le document, texte et figures compris.
% L'option « table » charge colortbl : elle donne \arrayrulecolor (filets
% colorés) et \rowcolors (fonds alternés) pour les tableaux de référence.
\usepackage[table]{xcolor}
\definecolor{inptblue}{HTML}{0082C8}   % accent : liens, numéros, mise en avant
\definecolor{navy}{HTML}{0B2E4F}       % titres
\definecolor{anrtred}{HTML}{BF162B}    % échec, abstention, danger
\definecolor{leafgreen}{HTML}{0E8F5E}  % modèle, succès
\definecolor{amberdark}{HTML}{B26A00}  % décision, garde-fou
\definecolor{plum}{HTML}{8B4FB0}       % quatrième série des figures de données
\definecolor{ink}{HTML}{1A1A1A}        % corps de texte
\definecolor{slate}{HTML}{5A6672}      % texte secondaire
\definecolor{mist}{HTML}{EEF2F6}       % fonds
\definecolor{rule}{HTML}{C9D3DC}       % filets

\color{ink}

% ---------------------------------------------------------------- géométrie
\usepackage[a4paper,left=3cm,right=2.5cm,top=2.7cm,bottom=2.7cm,
            headsep=14pt,footskip=28pt]{geometry}

\usepackage{setspace}
\setstretch{1.2}

% --------------------------------------------------------------- structures
\usepackage{graphicx}
\usepackage{float}
\usepackage{booktabs}
\usepackage{tabularx}
\usepackage{longtable}
\usepackage{multirow}
\usepackage{array}
\usepackage{enumitem}
\usepackage{afterpage}
\usepackage{emptypage}

\setlist{itemsep=2pt,parsep=0pt,topsep=4pt}
\setlist[itemize]{label=\textcolor{inptblue}{\small$\bullet$}}
\setlist[enumerate]{label=\textcolor{inptblue}{\bfseries\arabic*.}}

% Colonne X centrée / alignée à droite pour tabularx
\newcolumntype{C}{>{\centering\arraybackslash}X}
\newcolumntype{R}{>{\raggedright\arraybackslash}X}
\renewcommand{\arraystretch}{1.25}

% ------------------------------------------------------------------- titres
\usepackage{titlesec}

% Chapitre : numéro en très grand corps clair, titre en linéale foncée,
% filet fin de séparation. Pas de page de garde de chapitre.
\titleformat{\chapter}[display]
  {\normalfont\sffamily\bfseries\color{navy}}
  {\hfill\fontsize{72}{72}\selectfont\color{inptblue!35}\thechapter}
  {-14pt}
  {\Huge\raggedright}
  [\vspace{6pt}{\color{rule}\titlerule[1pt]}]

\titleformat{name=\chapter,numberless}[display]
  {\normalfont\sffamily\bfseries\color{navy}}
  {}
  {0pt}
  {\Huge\raggedright}
  [\vspace{6pt}{\color{rule}\titlerule[1pt]}]

\titlespacing*{\chapter}{0pt}{-24pt}{28pt}
\titlespacing*{name=\chapter,numberless}{0pt}{0pt}{28pt}

\titleformat{\section}
  {\normalfont\sffamily\Large\bfseries\color{navy}}
  {\color{inptblue}\thesection}{0.7em}{}
\titleformat{\subsection}
  {\normalfont\sffamily\large\bfseries\color{navy!90}}
  {\color{inptblue}\thesubsection}{0.6em}{}
\titleformat{\subsubsection}
  {\normalfont\sffamily\normalsize\bfseries\color{navy!80}}
  {\color{inptblue}\thesubsubsection}{0.5em}{}

\titlespacing*{\section}{0pt}{18pt}{8pt}
\titlespacing*{\subsection}{0pt}{14pt}{6pt}
\titlespacing*{\subsubsection}{0pt}{12pt}{4pt}

% Le paragraphe nommé sert de micro-titre dans les sections d'analyse.
\titleformat{\paragraph}[runin]
  {\normalfont\sffamily\bfseries\color{navy!85}}{}{0pt}{}[\;---]
\titlespacing*{\paragraph}{0pt}{10pt}{6pt}

\setcounter{secnumdepth}{3}
\setcounter{tocdepth}{2}

% ------------------------------------------------------------ en-têtes/pieds
\usepackage{fancyhdr}
\pagestyle{fancy}
\fancyhf{}
\renewcommand{\headrulewidth}{0pt}
\renewcommand{\footrulewidth}{0pt}
\fancyhead[L]{\footnotesize\sffamily\color{slate}\nouppercase{\leftmark}}
\fancyhead[R]{\footnotesize\sffamily\color{slate}\thepage}
% Filet sous l'en-tête, dessiné à la main pour en contrôler la couleur.
\renewcommand{\headrule}{\color{rule}\rule{\headwidth}{0.4pt}}
\renewcommand{\headrulewidth}{0.4pt}

\fancypagestyle{plain}{%
  \fancyhf{}%
  \renewcommand{\headrulewidth}{0pt}%
  \fancyfoot[C]{\footnotesize\sffamily\color{slate}\thepage}%
}

\renewcommand{\chaptermark}[1]{\markboth{\thechapter.\ #1}{}}

% ---------------------------------------------------------------- légendes
\usepackage[font=small,labelfont={sf,bf,color=navy},
            textfont={color=slate},labelsep=period,
            width=0.92\textwidth]{caption}
\captionsetup[figure]{skip=8pt}
\captionsetup[table]{skip=6pt,position=top}

% ------------------------------------------------------------------- encadrés
\usepackage[most]{tcolorbox}

% Encadré neutre : exemple, illustration.
\newtcolorbox{encadre}[1][]{%
  enhanced, breakable,
  colback=mist, colframe=inptblue,
  boxrule=0pt, leftrule=2.5pt,
  arc=0pt, outer arc=0pt,
  left=10pt, right=10pt, top=8pt, bottom=8pt,
  fonttitle=\sffamily\bfseries\color{navy},
  coltitle=navy, colbacktitle=mist,
  attach boxed title to top left={yshift=-2pt},
  boxed title style={boxrule=0pt, arc=0pt, left=0pt, right=0pt},
  #1}

% Encadré « constat » : un résultat que le lecteur doit retenir.
\newtcolorbox{constat}[1][]{%
  enhanced, breakable,
  colback=leafgreen!6, colframe=leafgreen,
  boxrule=0pt, leftrule=2.5pt,
  arc=0pt, outer arc=0pt,
  left=10pt, right=10pt, top=8pt, bottom=8pt,
  #1}

% Encadré « limite » : une réserve, une défaillance, un coût assumé.
\newtcolorbox{limite}[1][]{%
  enhanced, breakable,
  colback=anrtred!5, colframe=anrtred,
  boxrule=0pt, leftrule=2.5pt,
  arc=0pt, outer arc=0pt,
  left=10pt, right=10pt, top=8pt, bottom=8pt,
  #1}

% ---------------------------------------------------------------------- code
\usepackage{listings}
\lstdefinestyle{clair}{
  basicstyle=\ttfamily\footnotesize\color{ink},
  keywordstyle=\color{navy}\bfseries,
  commentstyle=\color{slate}\itshape,
  stringstyle=\color{leafgreen},
  numberstyle=\ttfamily\scriptsize\color{rule},
  backgroundcolor=\color{mist},
  frame=leftline, framerule=2pt, rulecolor=\color{rule},
  framexleftmargin=6pt, framexrightmargin=4pt,
  xleftmargin=8pt, xrightmargin=2pt,
  aboveskip=12pt, belowskip=8pt,
  breaklines=true, breakatwhitespace=false,
  columns=fullflexible, keepspaces=true,
  showstringspaces=false, tabsize=2,
  captionpos=b,
  literate={é}{{\'e}}1 {è}{{\`e}}1 {ê}{{\^e}}1 {à}{{\`a}}1
           {ç}{{\c c}}1 {û}{{\^u}}1 {î}{{\^\i}}1 {ô}{{\^o}}1
           {É}{{\'E}}1 {È}{{\`E}}1 {«}{{\guillemotleft}}1 {»}{{\guillemotright}}1
}
\lstset{style=clair}
\renewcommand{\lstlistingname}{Extrait}

% ------------------------------------------------------------------- figures
\usepackage{tikz}
\usepackage{pgfplots}
\usepackage{pgfplotstable}
\pgfplotsset{compat=1.18}
\usetikzlibrary{calc,positioning,arrows.meta,shapes.geometric,shapes.multipart,
                fit,backgrounds,decorations.pathreplacing,patterns,babel,
                chains,shadows.blur}
\usepgfplotslibrary{statistics}

%%%%%%%%%%%%%%%%%%%%%%%%%%%%%%%%%%%%%%%%%%%%%%%%%%%%%%%%%%%%%%%%%%%%%%%%%%%%%
%  Langage graphique commun à toutes les figures.
%
%  Les treize figures du rapport partagent ces styles : une forme et une
%  couleur y ont partout le même sens. Le lecteur n'a donc à apprendre la
%  légende qu'une seule fois.
%
%     gris    donnée, fichier, artefact          rectangle
%     bleu    traitement déterministe (code)     rectangle
%     vert    modèle appris (embeddings, LLM)    rectangle
%     ambre   décision, garde-fou                losange
%     rouge   arrêt, abstention, échec           rectangle
%
%  Palette de données validée pour la vision des couleurs (deutan/protan/
%  tritan) : bleu, ambre, vert, prune. Le rouge est réservé au statut
%  « échec » et n'est jamais employé comme simple série.
%%%%%%%%%%%%%%%%%%%%%%%%%%%%%%%%%%%%%%%%%%%%%%%%%%%%%%%%%%%%%%%%%%%%%%%%%%%%%

% ============================================================== SCHÉMAS TikZ

\tikzset{
  % --- boîtes ---------------------------------------------------------------
  bloc/.style={
    draw, rounded corners=2.5pt, line width=0.7pt,
    minimum height=11mm, text width=23mm, align=center,
    inner sep=3pt, font=\sffamily\scriptsize},
  donnee/.style={bloc, draw=slate, fill=mist,          text=ink},
  traitement/.style={bloc, draw=inptblue, fill=inptblue!8,  text=navy},
  modele/.style={bloc, draw=leafgreen, fill=leafgreen!8, text=leafgreen!45!black},
  arret/.style={bloc, draw=anrtred, fill=anrtred!7,    text=anrtred!70!black},
  decision/.style={
    draw=amberdark, fill=amberdark!10, line width=0.7pt,
    diamond, aspect=2.1, align=center, inner sep=1pt,
    text width=17mm, font=\sffamily\scriptsize, text=amberdark!60!black},
  % variantes de gabarit
  large/.style={text width=30mm},
  etroit/.style={text width=18mm},
  haut/.style={minimum height=14mm},
  %
  % --- liens ----------------------------------------------------------------
  flux/.style={-{Stealth[length=2.2mm,width=1.6mm]}, line width=0.65pt, draw=slate},
  fluxfort/.style={flux, line width=1pt, draw=navy},
  usage/.style={-{Stealth[length=2.2mm,width=1.6mm]}, line width=0.55pt,
                draw=slate!70, dash pattern=on 2pt off 1.6pt},
  rejet/.style={-{Stealth[length=2.2mm,width=1.6mm]}, line width=0.65pt, draw=anrtred},
  %
  % --- textes ---------------------------------------------------------------
  etiq/.style={font=\sffamily\scriptsize, text=slate, inner sep=2pt},
  etiqfort/.style={font=\sffamily\scriptsize\bfseries, text=navy, inner sep=2pt},
  annot/.style={font=\sffamily\scriptsize\itshape, text=slate, align=left,
                inner sep=2pt},
  %
  % --- conteneurs -----------------------------------------------------------
  bande/.style={
    draw=rule, line width=0.7pt, rounded corners=3pt,
    dash pattern=on 3pt off 2pt, inner sep=6mm},
  bandepleine/.style={draw=rule, line width=0.7pt, rounded corners=3pt,
                      fill=mist!45, inner sep=6mm},
  titrebande/.style={
    font=\sffamily\scriptsize\bfseries, text=slate,
    anchor=west, fill=white, inner xsep=4pt, inner ysep=1pt},
  %
  % --- pastille de renvoi vers le texte -------------------------------------
  puce/.style={
    circle, draw=anrtred, fill=white, text=anrtred, line width=0.6pt,
    inner sep=0.6pt, minimum size=3.6mm, font=\sffamily\bfseries\tiny},
}

% Titre posé sur le bord supérieur gauche d'un conteneur.
% #1 = nœud conteneur, #2 = libellé
%
% Le décalage passe par « xshift » et non par la syntaxe calc « $...$ » :
% celle-ci ne survit pas à un \resizebox, qui réinterprète le signe dollar.
\newcommand{\bandetitre}[2]{%
  \node[titrebande, xshift=6mm] at (#1.north west) {#2};}

% Légende compacte des types de nœud, à poser sous un schéma.
\newcommand{\legendeschema}{%
  \begin{tikzpicture}[baseline]
    \node[donnee, text width=0pt, minimum width=5mm, minimum height=3.2mm] (a) {};
    \node[etiq, right=1.2mm of a] (la) {donnée};
    \node[traitement, text width=0pt, minimum width=5mm, minimum height=3.2mm,
          right=4mm of la] (b) {};
    \node[etiq, right=1.2mm of b] (lb) {traitement};
    \node[modele, text width=0pt, minimum width=5mm, minimum height=3.2mm,
          right=4mm of lb] (c) {};
    \node[etiq, right=1.2mm of c] (lc) {modèle appris};
    \node[decision, text width=0pt, minimum width=5mm, minimum height=2.2mm,
          right=4mm of lc] (d) {};
    \node[etiq, right=1.2mm of d] (ld) {décision};
    \node[arret, text width=0pt, minimum width=5mm, minimum height=3.2mm,
          right=4mm of ld] (e) {};
    \node[etiq, right=1.2mm of e] {arrêt};
  \end{tikzpicture}}

% ============================================================ FIGURES pgfplots

% Séparateur décimal français dans toutes les figures : le corps du texte
% écrit « 0,969 », les graphiques doivent en faire autant.
\pgfkeys{/pgf/number format/.cd, use comma, 1000 sep={\,}}

\pgfplotsset{
  % Socle commun : pas de cadre, une seule grille (celle de l'axe des
  % valeurs), graduations discrètes. Tout ce qui n'aide pas à lire un chiffre
  % est retiré.
  socle/.style={
    width=\linewidth, height=6.4cm,
    font=\sffamily\small,
    axis line style={draw=rule, line width=0.6pt},
    axis x line*=bottom, axis y line*=left,
    every axis label/.style={font=\sffamily\small, text=slate},
    tick label style={font=\sffamily\scriptsize, text=slate},
    label style={font=\sffamily\scriptsize, text=slate},
    title style={font=\sffamily\small\bfseries, text=navy, align=left},
    major tick length=2.5pt, minor tick length=0pt,
    tick align=outside,
    grid style={draw=rule!55, line width=0.4pt},
    legend style={
      draw=none, fill=none, font=\sffamily\scriptsize, text=slate,
      cells={anchor=west}, legend columns=-1,
      /tikz/every even column/.append style={column sep=8pt}},
    legend image post style={line width=2pt},
    clip=false,
  },
  % Barres horizontales : la catégorie se lit à plat, sans tête penchée.
  barresH/.style={
    socle,
    xbar, bar width=9pt,
    xmajorgrids, ymajorgrids=false,
    ytick=data, y dir=reverse,
    enlarge y limits=0.14,
    % L'axe vertical ne porte aucune quantité : on efface son trait et ses
    % graduations, mais surtout PAS « axis y line=none », qui emporterait
    % aussi les libellés de catégorie.
    axis x line*=bottom, axis y line*=left,
    y axis line style={draw=none},
    ytick style={draw=none},
    nodes near coords style={font=\sffamily\scriptsize, text=ink,
                             anchor=west, xshift=2pt},
  },
  % Barres verticales groupées.
  barresV/.style={
    socle,
    ybar, bar width=13pt,
    ymajorgrids, xmajorgrids=false,
    enlarge x limits=0.22,
    xtick=data,
    nodes near coords style={font=\sffamily\scriptsize, text=ink,
                             anchor=south, yshift=1pt},
  },
  % Nuage de points en bandes (scores d'abstention).
  bandesY/.style={
    socle,
    ymajorgrids=false, xmajorgrids=true,
    axis y line=none,
    enlarge y limits=0.3,
  },
}

% Cycle de couleurs des séries — ordre fixe, jamais recyclé.
\pgfplotscreateplotcyclelist{serie}{
  {draw=inptblue,  fill=inptblue},
  {draw=amberdark, fill=amberdark},
  {draw=leafgreen, fill=leafgreen},
  {draw=plum,      fill=plum},
}
\pgfplotsset{cycle list name=serie}

% Extrémité de barre arrondie côté valeur, ancrée sur la ligne de base.
\pgfplotsset{
  bout arrondi/.style={
    /pgfplots/every axis plot post/.append style={
      rounded corners=1.6pt}},
}

% Filet vertical de seuil, avec son étiquette.
% #1 = abscisse, #2 = libellé
\newcommand{\seuil}[2]{%
  \draw[anrtred, line width=0.8pt, dash pattern=on 3pt off 2pt]
    ({axis cs:#1,0}|-{rel axis cs:0,0}) -- ({axis cs:#1,0}|-{rel axis cs:0,1});
  \node[font=\sffamily\scriptsize\bfseries, text=anrtred, anchor=south, rotate=90]
    at ({axis cs:#1,0}|-{rel axis cs:0,0.02}) {\;#2};}

% ================================================================ DIAGRAMME
% Gantt : mêmes couleurs, mêmes fontes que le reste.
\ganttset{
  bar height=0.55,
  bar/.append style={draw=inptblue!70, fill=inptblue!25, rounded corners=1pt},
  bar label font=\sffamily\scriptsize\color{ink},
  group/.append style={draw=navy, fill=navy},
  group label font=\sffamily\scriptsize\bfseries\color{navy},
  milestone/.append style={draw=amberdark, fill=amberdark},
  milestone label font=\sffamily\scriptsize\color{amberdark!60!black},
  title label font=\sffamily\scriptsize\color{slate},
  title/.append style={draw=none, fill=none},
  link/.style={-{Stealth[length=1.6mm]}, draw=slate, line width=0.5pt},
  vgrid={*1{draw=rule!60, line width=0.4pt}},
}


% ---------------------------------------------------------------- références
\usepackage{glossaries}

\usepackage[hidelinks]{hyperref}
\hypersetup{
  colorlinks=true,
  linkcolor=navy,
  citecolor=inptblue,
  urlcolor=inptblue,
  pdfauthor={Taha El Badaoui},
  pdftitle={Assistant juridique souverain — RAG sur le Code du travail marocain},
  pdfsubject={Rapport de stage de fin de première année — INPT, filière ASEDS},
  pdfkeywords={RAG, souveraineté numérique, TAL, recherche sémantique, droit du travail},
  pdfpagemode=UseOutlines,
  bookmarksnumbered=true,
}
\urlstyle{same}

% --------------------------------------------------------------- raccourcis
\newcommand{\module}[1]{\texttt{\small #1}}
\newcommand{\anglais}[1]{\textit{#1}}
\newcommand{\artref}[1]{article~#1}

% Page blanche technique (après la page de garde).
\newcommand{\pageblanche}{%
  \null\thispagestyle{empty}\addtocounter{page}{-1}\newpage}

% Chiffres générés depuis les fichiers de mesure — voir rapport/build_data.py.
% Fichier généré par rapport/build_data.py — NE PAS ÉDITER.
% Régénérer avec :  python rapport/build_data.py

\newcommand{\NbArticles}{589}
\newcommand{\NbLivres}{8}
\newcommand{\NbQuestions}{62}
\newcommand{\NbInScope}{53}
\newcommand{\NbOutScope}{9}
\newcommand{\NbRepondues}{51}
\newcommand{\NbRefusTort}{2}
\newcommand{\NbGarde}{2}
\newcommand{\TauxCitation}{96}
\newcommand{\TauxVerif}{81}
\newcommand{\TauxSourceAttendue}{88}
\newcommand{\TauxAbstention}{100}
\newcommand{\NbCitations}{95}
\newcommand{\NbNonVerifiees}{18}
\newcommand{\SeuilAbstention}{0{,}42}
\newcommand{\ScoreInMin}{0{,}596}
\newcommand{\ScoreInMax}{0{,}793}
\newcommand{\ScoreJuridiqueMin}{0{,}445}
\newcommand{\ScoreJuridiqueMax}{0{,}603}
\newcommand{\ScoreHorsDroitMin}{0{,}307}
\newcommand{\ScoreHorsDroitMax}{0{,}446}
\newcommand{\LatRecherche}{0{,}16}
\newcommand{\LatGeneration}{3{,}83}
\newcommand{\LatTotale}{3{,}98}
\newcommand{\LatAbstention}{0{,}13}
\newcommand{\LatMax}{16{,}32}
\newcommand{\PartRecherche}{4}
\newcommand{\FacteurGarde}{30}
\newcommand{\RecallDense}{0{,}849}
\newcommand{\MrrDense}{0{,}778}
\newcommand{\RecallBmXXV}{0{,}660}
\newcommand{\MrrBmXXV}{0{,}535}
\newcommand{\RecallHybride}{0{,}868}
\newcommand{\MrrHybride}{0{,}775}
\newcommand{\ModeleLLM}{mistral:\allowbreak 7b}
\newcommand{\DateMesure}{2026-07-31}

